\section*{Discussion}

The measurements are consistent with a single mechanism. Parity labels make a symmetry-mandated zero structural: it holds at any parameter values, cannot be trained away because it was never trained in, and differentiates correctly with respect to atomic displacements. Removing the labels does not merely degrade the prediction --- it converts a certainty into an ordinary regression output, which on centrosymmetric crystals is an ordinary regression error of the magnitude of the property itself. The dipole control shows the failure is not a property of odd-parity targets as such but of evaluating them where symmetry makes the answer exact; the point-group split shows that what SO(3) models retain is precisely the rotational part of Neumann's principle and nothing more; and the agreement between architecturally unrelated models shows the effect belongs to the symmetry class, not to a particular network.

The result gives quantitative content to a distinction the field usually treats as terminology. Curie's principle --- that an equivariant map cannot produce outputs less symmetric than its inputs --- is the foundation on which equivariant networks are advertised, and an active literature engineers controlled ways to relax it when physical symmetry breaking demands. Our measurements show the principle is silently violated, for improper operations, by a widely deployed subclass of the very models the principle is invoked to describe. For tensor properties, upholding Curie's principle requires equivariance to the full orthogonal group; rotational equivariance alone upholds only its proper half.

The distinction also matters for how symmetry-correct predictions should be engineered. Space-group-informed enforcement modules, canonicalisation schemes, and symmetry-adapted output decompositions repair the output of a model whose features remain parity-blind, and our test-time symmetrisation experiment makes the cost of that arrangement concrete: the repair works only for the quantity someone remembered to repair, requires knowing the target's parity in advance, and leaves every other quantity derived from the same features unprotected. Feature-level parity, by contrast, protects everything the features feed --- including derivatives, as the Jacobian analysis shows. This is not an argument against readout-level methods, which have accuracy benefits of their own; it is an argument that the feature symmetry group is a prior and independent design decision, made once for the whole model, and currently made invisibly. As property heads are increasingly attached to embeddings of pretrained universal potentials, that decision is inherited by every downstream task: a parity-even rank-2 target such as the dielectric tensor never exposes it, and the first parity-odd head attached to a rotation-only backbone will fail exactly as measured here.

An apparent tension with experiment deserves a comment. Sensitive resonance techniques detect piezoelectric responses in nominally centrosymmetric materials\cite{aktas2021nominally}, and one might ask whether an SO(3) model's nonzero predictions anticipate such physics. They do not, in two ways. The measured symmetry-forbidden responses arise because the actual structure of a real sample deviates from its nominal average symmetry --- defects, disorder, incomplete transitions --- and their reported effective coefficients fall orders of magnitude below the responses of conventional piezoelectrics, whereas the SO(3) violations here sit on exactly symmetric digital structures and reach the full magnitude of genuine piezoelectric tensors. The parity-labelled models, meanwhile, capture the legitimate version of the phenomenon: the single raw-coordinate structure an O(3) arm flagged turned out to be genuinely non-centrosymmetric below the selection tolerance, meaning the model responded to real symmetry breaking present in the coordinates, which is the same physics the resonance experiments exploit.

EquiformerV2's behaviour is reported strictly as measured properties of its released code, reproduced against the pinned upstream source so that none of it is an artifact of our packaging: its evaluation draws a random per-edge reference frame on every forward pass and is therefore nondeterministic, its rotational equivariance is approximate with relative errors of order ten percent, and its autograd Jacobian is incomplete because the Wigner matrices are detached from atomic positions (Supplementary Note~10 and Supplementary Tables~15--16). None of our conclusions depends on these properties; we note that for an exactly rotation-equivariant model the output would be independent of the arbitrary frame, so the observed evaluation spread is itself a direct readout of equivariance error. Its inclusion serves one purpose: it shows the headline behaviour in a deployed model that was never touched by our matched-pair construction.

The study has limits. All three matched O(3) arms realise their parity labels through the same irreducible-representation library, e3nn\cite{geiger2022e3nn}; an intended non-e3nn O(3) control was withdrawn for numerical conditioning reasons unrelated to parity, an instructive failure documented in Supplementary Note~11, so we cannot fully exclude implementation-specific contributions to the O(3) side. Three seeds per configuration limit seed-level statistics, and all statistical weight comes from the structure-level paired tests. The violation magnitude is extensive, so absolute values depend on system size, although a size-normalised variant of the headline metric leaves every conclusion unchanged (Supplementary Note~3). The scalar-energy control is far from converged and its absolute accuracy is not representative of tuned models, though both arms share its configuration exactly. The augmentation experiment used one budget and schedule, and the loss-weight sweep one architecture; neither excludes that more aggressive regimes could push violations lower, but the in-training control bounds what any such regime can achieve --- approaching zero is not attaining it, and the physics requires exactness. Finally, hyperparameters were shared across cores rather than tuned per architecture, which is the price of the matched-pair design.

The practical conclusion is compact. Whether a machine-learned model can make physically impossible tensor predictions is decided before training begins, by whether its features carry parity labels --- a property that most released architectures lack and none advertises. For parity-odd targets, and for any pipeline that may one day attach one, the label is not bookkeeping; it is the physics.
